\subsection{Contextos matemáticos y preprocesamiento}
\label{sec:dl_preproc}

El análisis de señales acústicas mediante \emph{Deep Learning} requiere transformar el audio crudo en representaciones que capturen adecuadamente sus propiedades espectrales y temporales.  
Para ello se emplean herramientas de la teoría de Fourier y representaciones perceptuales que aproximan el comportamiento del sistema auditivo humano.  
En lo que sigue se presentan las bases matemáticas de tres transformaciones fundamentales: la Transformada Rápida de Fourier (FFT), la Transformada de Fourier de Tiempo Corto (STFT) y la escala Mel.

\subsubsection{Transformada Rápida de Fourier (FFT)}

Dada una señal discreta $x[n]$ de longitud $N$, su Transformada Discreta de Fourier (DFT) está definida como
\begin{equation}
	X[k] = \sum_{n=0}^{N-1} x[n] \, e^{-j \tfrac{2\pi}{N} kn}, 
	\quad k = 0,1,\dots,N-1.
	\label{eq:dft}
\end{equation}

La DFT descompone la señal en senoidales complejas ortogonales, revelando la distribución de energía en frecuencia \cite{Haykin2001,Proakis2007}.  
Su cálculo directo tiene costo $\mathcal{O}(N^2)$, pero el algoritmo FFT lo reduce a $\mathcal{O}(N \log N)$, habilitando aplicaciones en tiempo real como detección acústica.  

Es importante notar que la DFT brinda una descripción global: indica magnitud y fase de cada componente frecuencial, pero pierde información sobre el instante en que ocurrió.  
En bioacústica esto es relevante: el zumbido de un mosquito puede variar en fracciones de segundo, lo que motiva el uso de representaciones con resolución conjunta en tiempo y frecuencia.

\insertarfigura{fig/fft_signal.png}
{Conversión de una señal temporal a su representación espectral mediante FFT.}
{fig:fft_signal}{1}

\subsubsection{Transformada de Fourier de Tiempo Corto (STFT)}

La STFT extiende la DFT incorporando localización temporal.  
La señal se multiplica por una ventana $w[n]$ de duración finita, que selecciona un segmento, y se desplaza en el tiempo mediante un parámetro $m$. Su definición es:
\begin{equation}
	X(m,\omega) = \sum_{n=-\infty}^{\infty} x[n] \, w[n-m] \, e^{-j \omega n}.
	\label{eq:stft}
\end{equation}

El resultado es una función bidimensional que describe cómo varía el contenido espectral en el tiempo, dando lugar al \emph{espectrograma}.  
La resolución temporal y frecuencial dependen inversamente de la longitud de la ventana:  
ventanas cortas capturan sucesos breves pero con poca resolución espectral, mientras que ventanas largas ofrecen buena resolución en frecuencia pero pierden detalle temporal.  
Esta dualidad, conocida como principio de incertidumbre tiempo–frecuencia, condiciona el diseño del preprocesamiento \cite{Bancroft2002}.  

En detección de mosquitos, la STFT permite observar la evolución dinámica del zumbido, que puede variar en frecuencia fundamental por el movimiento del insecto o por interferencias ambientales.

\insertarfigura{fig/stft_spectrogram.png}
{Esquema de la STFT mostrando segmentación por ventanas.}
{fig:stft_spectrogram}{1}

\subsubsection{Escala Mel}

Aunque la STFT proporciona un espectro en Hz, el sistema auditivo humano no percibe las frecuencias de manera lineal.  
Somos más sensibles a bajas frecuencias: diferenciar entre 100 y 200 Hz resulta más fácil que entre 5000 y 5100 Hz.  
Para aproximar esta percepción se usa la escala Mel, que relaciona la frecuencia física $f$ con una perceptual $f_{\text{mel}}$:
\begin{equation}
	f_{\text{mel}} = 2595 \, \log_{10}\!\left( 1 + \frac{f}{700} \right).
	\label{eq:mel}
\end{equation}

El cálculo de un \emph{mel-spectrograma} consiste en aplicar un banco de filtros triangulares, distribuidos equidistantemente en la escala Mel, sobre el espectrograma lineal.  
La energía espectral se proyecta a un dominio perceptual comprimido, reduciendo dimensionalidad e introduciendo cierta invariancia frente a variaciones irrelevantes para el oído humano.  

En \emph{Deep Learning} aplicado a audio, el mel-spectrograma es una de las representaciones más utilizadas como entrada a CNNs, pues concentra la información discriminativa y mejora la robustez frente a ruido ambiental \cite{Campos2022,HumBugDB}.

\insertarfigura{fig/mel_spectrogram.png}
{Waveform, espectrograma y mel-spectrograma de un segmento de audio de 10 segundos (Google AudioSet). 
	El mel-spectrograma, basado en la escala perceptual Mel, ofrece mayor resolución en bajas frecuencias que el espectrograma lineal.}
{fig:mel_spectrogram}{0.9}
